\title{Controlling Text-to-Image Diffusion by Orthogonal Finetuning}

\begin{abstract}
Text-to-image diffusion models have recently shown remarkable proficiency in synthesizing lifelike images from textual descriptions. However, guiding or steering these robust models to accomplish distinct downstream applications remains an outstanding and significant research issue. In response, we propose Orthogonal Finetuning~(OFT), a theoretically grounded adaptation approach for refining text-to-image diffusion models toward specific downstream objectives. In contrast with existing approaches, OFT is theoretically guaranteed to conserve the hyperspherical energy, a metric quantifying pairwise neuron relationships constrained on the unit hypersphere. Our investigations demonstrate that maintaining this property is essential for upholding the semantic image generation capabilities inherent to text-to-image diffusion models. To further enhance the reliability of finetuning, we introduce Constrained Orthogonal Finetuning (COFT), which adds an extra constraint enforcing a fixed radius on the hypersphere. Our work focuses on two central text-to-image finetuning challenges: subject-driven generation—producing images depicting a specific subject in diverse scenarios with a few exemplars and a prompt; and controllable generation—enabling the model to incorporate auxiliary control inputs. Through empirical studies, we show that the OFT framework delivers superior generation quality and accelerates convergence in comparison with existing methods.
\end{abstract}

\section{Introduction}

State-of-the-art text-to-image diffusion models~\cite{saharia2022photorealistic,ramesh2022hierarchical,rombach2022high} have achieved remarkable advances in producing visually plausible images conditioned on textual cues. Nevertheless, textual input alone may often lack sufficient specificity and precision to impart detailed and reliable control over the resulting imagery. To mitigate these limitations, this paper investigates two crucial text-to-image generation paradigms:

\begin{itemize}[leftmargin=*,nosep]

    \item \textbf{Subject-driven generation}~\cite{ruiz2023dreambooth}: The objective is to create novel images portraying the same entity in various contexts, based on a limited collection of the subject’s images and guidance from a text prompt.
    \item \textbf{Controllable generation}~\cite{zhang2023adding,mou2023t2i}: In this scenario, an additional control input (such as canny edge maps or segmentation masks) is provided to direct the image synthesis process in tandem with textual guidance.
\end{itemize}

At their core, both tasks address the challenge of how to adapt text-to-image diffusion models to new tasks through finetuning, while maintaining the generative strengths acquired during pretraining. An ideal finetuning method should fulfill the following criteria: (1) \emph{training efficiency}, meaning that it involves a minimal number of trainable parameters and requires fewer training epochs; and (2) \emph{generalizability preservation}, that is, it should retain the model’s ability to generate high-quality and varied outputs. Common approaches for finetuning involve either subtly modifying the model’s neuron weights using a small learning rate (\emph{e.g.}, \cite{ruiz2023dreambooth}), or augmenting the network with a compact module whose parameters are adjusted (\emph{e.g.}, \cite{hulora2022,zhang2023adding}). Although straightforward, these methods do not inherently safeguard the original generative capabilities conferred by pretraining. Furthermore, there is a notable gap in systematic principles for crafting robust finetuning strategies or for selecting appropriate hyperparameters such as training epochs. The primary challenge stems from the absence of a quantitative metric for assessing how well the original generative abilities are preserved. Most current finetuning techniques operate under the assumption that maintaining a small Euclidean distance between the parameters of the finetuned model and those of the pretrained model will preserve generative performance, which motivates the use of very small learning rates. However, while this may sometimes suffice, we contend that Euclidean distance alone is insufficient for encapsulating semantic fidelity. Thus, adopting a more structurally informed metric to compare finetuned and pretrained models could substantially enhance both the preservation of generative performance and the stability of the finetuning process.

Building on prior empirical findings that hyperspherical similarity effectively captures semantic structure~\cite{liu2017deep,liu2018decoupled,chen2020angular}, we leverage hyperspherical energy~\cite{liu2018learning} as a means to describe the inter-neuron relationships within a layer. This metric, computed as the aggregate hyperspherical similarity (such as cosine similarity) across every neuron pair in a given layer, quantifies how uniformly neurons are distributed over the unit hypersphere~\cite{liu2021learning}. We posit that, after finetuning, a model should exhibit a hyperspherical energy closely matching that of its pretrained counterpart. One straightforward strategy would be to introduce a regularization term that preserves hyperspherical energy during finetuning; however, this approach lacks assurances of effectively minimizing the energy deviation. Alternatively, we exploit a key invariance property: pairwise hyperspherical similarities remain intact under any common orthogonal transformation applied to all neurons. This observation leads us to Orthogonal Finetuning~(OFT), a technique for adapting large-scale text-to-image diffusion models to new tasks without altering their hyperspherical energy. OFT operates by learning a shared orthogonal transformation per layer to maintain the angular relationships among neurons, effectively realigning the canonical basis for each layer's neurons. Additionally, recognizing that maintaining a smaller Euclidean distance between finetuned and pretrained models helps to conserve original performance, we introduce a refinement—Constrained Orthogonal Finetuning~(COFT)—which restricts finetuned parameters to lie on a hypersphere with fixed radius centered at the pretrained neuron positions.

The rationale behind employing orthogonal transformation for finetuning neurons is partially inspired by the principles of the 2D Fourier transform, where an image can be decomposed into its magnitude and phase spectra. The phase spectrum, reflecting the angular relation between the input and basis, is largely responsible for conveying semantic information. Notably, if the phase spectrum of an image is retained while substituting a random magnitude spectrum, the reconstructed image still maintains most of its semantic content. This observation highlights the importance of altering neuron orientations for semantic modification of generated images—a central objective in both subject-driven and controllable generation tasks. Nevertheless, allowing neuron orientations to vary with excessive freedom may compromise the generative capabilities established during pretraining. To address this, we propose to keep the inter-neuronal angles invariant, based on the premise that these angles are integral for encoding knowledge within neural networks. Consequently, we adopt a layer-wise, shared orthogonal transformation for neurons within each layer, ensuring that the hyperspherical energy remains consistent.

Our approach is also motivated by prior work on orthogonal over-parameterized training~\cite{liu2021orthogonal}, which involves constructing classification neural networks from the ground up by applying orthogonal transformations to a randomly initialized network. A randomly initialized network is expected, on average, to exhibit low hyperspherical energy, and the intent in~\cite{liu2021orthogonal} is to maintain this low energy state throughout training (since reduced hyperspherical energy correlates with enhanced generalization in classification tasks~\cite{liu2018learning,lin2020regularizing}). It has been demonstrated that orthogonal transformations offer enough expressivity to train well-generalizing classifiers. In our work, the focus shifts to finetuning text-to-image diffusion models, targeting enhanced controllability and superior generative quality for downstream tasks. The distinction between our Orthogonal Finetuning Transformation (OFT) and~\cite{liu2021orthogonal} is twofold. First, unlike~\cite{liu2021orthogonal}, which actively minimizes hyperspherical energy, OFT aims to preserve the energy level of the pretrained model, thus safeguarding its inherent structure during finetuning; reducing the energy in this context may disrupt essential semantic representations. Second, OFT explicitly constrains deviations from the pretrained configuration, leading to a variant with such restrictions, while~\cite{liu2021orthogonal} does not introduce comparable constraints. The central challenge in finetuning lies in finding an appropriate balance between adaptability and model integrity. We argue that the OFT methodology effectively navigates this trade-off. Our key contributions are summarized as follows:

\begin{itemize}[leftmargin=*,nosep]

    \item We introduce a new finetuning approach termed Orthogonal Finetuning (OFT), designed to enhance the controllability of text-to-image diffusion models. Additionally, we develop a constrained extension that restricts angular deviations from the original pretrained model, thereby bolstering stability.
    \item Unlike prevailing finetuning strategies, OFT modifies model parameters while mathematically guaranteeing the conservation of hyperspherical energy—a property we empirically identify as critical for maintaining the pretrained model's generative semantic integrity.
    \item We implement OFT for both subject-driven and controllable generation tasks. Through thorough experimentation, we observe marked advancements over prior techniques in generation fidelity, training convergence speed, and overall finetuning robustness. OFT is also shown to be more sample-efficient, converging successfully with far fewer training images and epochs.
    \item In the context of controllable generation, we propose a novel control consistency metric for assessing controllability. This metric functions by inferring the control signal from the output image and measuring its similarity to the original control input, providing further empirical support for the efficacy of OFT.
\end{itemize}

\section{Related Work}

\textbf{Text-to-image diffusion models}. Recent years have witnessed substantial breakthroughs~\cite{saharia2022photorealistic,ramesh2022hierarchical,rombach2022high,gu2022vector,nichol2022glide} in text-to-image generation, propelled by advancements in diffusion-based generative modeling~\cite{ho2020denoising,song2021score,song2021denoising,dhariwal2021diffusion} and the evolution of vision-language representation techniques~\cite{su2020vlbert,lu2019vilbert,radford2021learning,wang2021simvlm,li2022blip,li2023blip,alayrac2022flamingo,schuhmann2022laion}. Approaches such as GLIDE~\cite{nichol2022glide} and Imagen~\cite{saharia2022photorealistic} focus on pixel-space diffusion models; GLIDE involves joint training of a text encoder and diffusion prior using paired modalities, while Imagen leverages a frozen, pretrained text encoder. Conversely, models like Stable Diffusion~\cite{rombach2022high} and DALL-E2~\cite{ramesh2022hierarchical} employ latent space diffusion modeling. Stable Diffusion’s latent space is structured via a visual codebook learned by VQ-GAN~\cite{esser2021taming}, while DALL-E2 constructs a joint text-image embedding with CLIP~\cite{radford2021learning}. Beyond diffusion approaches, text-to-image synthesis has also been advanced through generative adversarial networks~\cite{reed2016generative,zhang2017stackgan,xu2018attngan,li2019controllable} and autoregressive models~\cite{ramesh2021zero,ding2021cogview,wu2022nuwa,yuscaling}. Notably, OFT is a model-agnostic finetuning scheme suitable for application across various text-to-image diffusion architectures.

\textbf{Subject-driven generation}. Strategies to preserve subject characteristics include utilizing user-provided masks as supplementary conditioning, as explored in \cite{avrahami2022blended,nichol2022glide}. Inversion-based techniques~\cite{choi2021ilvr,dhariwal2021diffusion,ramesh2022hierarchical,gal2022image} facilitate context modifications while maintaining subject consistency, and approaches like \cite{hertz2022prompt} permit both local and global edits absent explicit masks. Nonetheless, such methods frequently fall short in retaining intricate, identity-specific subject details. Pivotal Tuning~\cite{roich2022pivotal} addresses this by finetuning the generator near an inverted latent vector, incorporating regularization for identity retention, while \cite{nitzan2022mystyle} derives a personalized face prior from a set of images of an individual. \cite{casanova2021instance} enables varied instantiations of a given instance, yet often sacrifices fidelity to the original subject. DreamBooth~\cite{ruiz2023dreambooth} introduces a customized token and finetunes the diffusion model on a small set of subject images using a blend of reconstruction and prior preservation losses. Building upon the DreamBooth methodology, OFT diverges by replacing straightforward low-rate finetuning with a structured orthogonal transformation-based update to the model.

\textbf{Controllable generation}. Image-to-image translation represents a subset of controllable generation, with earlier approaches largely leveraging conditional generative adversarial networks~\cite{isola2017image,zhu2017toward,wang2018high,park2019semantic,choi2018stargan}. More recently, diffusion models have also been adapted for image-to-image translation tasks~\cite{saharia2022palette,voynov2022sketch,wang2022pretraining}. ControlNet~\cite{zhang2023adding} advances controllability by adapting a pretrained diffusion model through finetuning in response to supplementary control signals, achieving notable results in this area. In parallel, T2I-Adapter~\cite{mou2023t2i} enhances the controllability of generated outputs by finetuning pretrained diffusion models. Aligning with the experimental design in \cite{zhang2023adding,mou2023t2i}, we employ OFT to finetune pretrained diffusion models. This approach consistently offers enhanced controllability with a reduced need for training samples and finetuning parameters. Importantly, OFT avoids introducing any extra computational cost during inference.

\textbf{Model finetuning}. The practice of finetuning large-scale pretrained models on specialized downstream tasks has become standard~\cite{devlin2018bert,brown2020language,he2022masked}. Adaptation-based finetuning approaches (\emph{e.g.}, \cite{hulora2022,houlsby2019parameter,pfeiffer2020adapterfusion}) are extensively explored within natural language processing. OFT is closely related to LoRA~\cite{hulora2022}, which models additive parameter updates with a low-rank structure during finetuning. However, OFT distinguishes itself by applying a shared orthogonal transformation across layers to multiplicatively modify neuron weights, which guarantees the preservation of pairwise angular relationships between neurons in a layer. This results in greater model stability.

\section{Orthogonal Finetuning}

\subsection{Why Does Orthogonal Transformation Make Sense?}

We begin by motivating the suitability of orthogonal transformations in finetuning text-to-image diffusion models. This rationale is unpacked by addressing two subsidiary considerations: (1) the importance of finetuning neuron angles (\emph{i.e.}, their directional orientation), and (2) the reasoning for utilizing orthogonal transformations as the mechanism for angle adjustment.

Addressing the first question, we take cues from empirical findings reported in \cite{liu2018decoupled,chen2020angular}, which suggest that differences in angular features effectively reflect semantic disparities. SphereNet~\cite{liu2017deep} further demonstrates that constraining all neuron outputs to lie on a unit hypersphere preserves network expressiveness and can enhance generalization performance. This observation highlights that the orientation of neurons—rather than their magnitude—encapsulates crucial information from the data. To clarify the influence of neuron angles, we present a toy experiment depicted in Figure~\ref{fig:toy_ae}. Here, a typical convolutional autoencoder is trained on a dataset of flower images. During training, we employ a standard inner product to compute the feature map, where $z$ denotes the convolution kernel’s response (with $\bm{w}$ as kernel weights and $\bm{x}$ as the windowed input). At the evaluation stage, three distinct approaches to feature map generation are examined: (a) repeating the inner product used during training, (b) considering only the response magnitude, and (c) using solely angular (directional) information. As illustrated in Figure~\ref{fig:toy_ae}, the angular component alone is capable of nearly perfect image reconstruction, whereas magnitude alone fails to convey meaningful information. Importantly, cosine-based angular activations (c) are not utilized during training; all learning occurs through inner products. These results point to the pivotal role of neuron directionality in encoding the semantic content of input images. Hence, adjusting neuron directions appears particularly promising for semantic modification tasks.

Turning to the second question, the most direct approach to refining neuron directions is by collectively rotating or reflecting the neurons within the same layer, which inherently involves applying an orthogonal transformation. Although more flexible angular adjustments could involve rotating each neuron individually by different amounts, our findings indicate that orthogonal transformations achieve a favorable balance between flexibility and structural regularization. Supporting this, \cite{liu2021orthogonal} attests to the capacity of orthogonal transformations to facilitate effective neural network learning. To empirically validate this, we run a controlled generation experiment within the ControlNet~\cite{zhang2023adding} setup, presenting the findings in Figure~\ref{fig:ortho}. The results show that our standard implementation of OFT delivers stable performance and precise control post-training (at epoch 20). In contrast, when the orthogonality constraint is lifted, OFT fails to yield convincing outputs or controllability. This underscores the essential contribution of orthogonality constraints in the success of OFT.

\subsection{General Framework}

In standard finetuning approaches, models—or subsets of their layers—are updated via gradient descent, typically employing a small learning rate. This conservative learning rate inherently restricts the deviation from the pretrained model parameters, with classical finetuning procedures effectively optimizing the objective $\thickmuskip=2mu \medmuskip=2mu \| \bm{M} - \bm{M}^0\|$, where $\bm{M}$ represents the adapted model weights, and $\bm{M}^0$ denotes those from pretraining. Although such an implicit constraint is reasonable, it may still offer excessive flexibility, particularly when adapting large-scale models. 

To mitigate this, LoRA augments the procedure by introducing a low-rank restriction on the weight modifications, formally written as $\thickmuskip=2mu \medmuskip=2mu \text{rank}(\bm{M} - \bm{M}^0)=r'$, with $r'$ being a small, fixed integer. In contrast, OFT implements a constraint targeting neuron-level relationships: it enforces $\thickmuskip=2mu \medmuskip=2mu \| \text{HE}(\bm{M})-\text{HE}(\bm{M}^0) \|=0$, in which $\text{HE}(\cdot)$ denotes the hyperspherical energy associated with the weight matrix. 

To illustrate, consider a fully connected layer $\thickmuskip=2mu \medmuskip=2mu \bm{W}=\{\bm{w}_1,\cdots,\bm{w}_n\}\in\mathbb{R}^{d\times n}$, with the $i$-th neuron vector given by $\bm{w}_i\in\mathbb{R}^{d}$ (and $\bm{W}^0$ its pretrained counterpart). The output $\thickmuskip=2mu \medmuskip=2mu \bm{z}\in\mathbb{R}^n$ is generated as $\thickmuskip=2mu \medmuskip=2mu \bm{z}=\bm{W}^\top\bm{x}$, where $\thickmuskip=2mu \medmuskip=2mu \bm{x}\in\mathbb{R}^d$ is the input feature vector. The OFT methodology can thus be conceptualized as minimizing the change in hyperspherical energy between pretrained and finetuned weights:

\begin{equation}\label{eq:oft_principle}
\footnotesize
\min_{\bm{W}} \left\| \text{HE}(\bm{W})-\text{HE}(\bm{W}^0)\right\|~~\Leftrightarrow~~\min_{\bm{W}} \bigg{\|} \sum_{i\neq j} \|\hat{\bm{w}}_i-\hat{\bm{w}}_j\|^{-1}-\sum_{i\neq j} \|\hat{\bm{w}}^0_i-\hat{\bm{w}}^0_j\|^{-1}\bigg{\|}
\end{equation}

Here, each normalized neuron is represented as $\thickmuskip=2mu \medmuskip=2mu \hat{\bm{w}}_i:=\bm{w}_i/\|\bm{w}_i\|$, and the hyperspherical energy for $\bm{W}$ is computed as $\thickmuskip=2mu \medmuskip=2mu \text{HE}(\bm{W}):= \sum_{i\neq j} \|\hat{\bm{w}}_i-\hat{\bm{w}}_j\|^{-1}$. The objective in Eq.~\eqref{eq:oft_principle} attains a minimum value of zero, a condition met whenever $\bm{W}$ differs from $\bm{W}^0$ solely by a rotation or reflection, that is, $\thickmuskip=2mu \medmuskip=2mu \bm{W}=\bm{R}\bm{W}^0$ with an orthogonal matrix $\thickmuskip=2mu \medmuskip=2mu \bm{R}\in\mathbb{R}^{d\times d}$ (with determinant $1$ for rotation or $-1$ for reflection). This observation forms the foundation of OFT, whereby fine-tuning is performed by learning orthogonal matrices (shared across layers) that transform the neurons in each layer. More formally, for a given pretrained fully connected layer

\begin{equation}\label{eq:oft_general}
    \footnotesize
    \bm{z}=\bm{W}^\top\bm{x}=(\bm{R}\cdot\bm{W}^0)^\top\bm{x},~~~\text{s.t.}~\bm{R}^\top\bm{R}=\bm{R}\bm{R}^\top=\bm{I}
\end{equation}

In this equation, $\bm{W}$ is the weight matrix obtained through OFT-finetuning, and $\bm{I}$ denotes the identity matrix. The OFT approach is depicted in Figure~\ref{fig:block}. Analogous to the use of zero initialization in LoRA, it is important for OFT to ensure that fine-tuning is performed beginning exactly from the pretrained parameter $\bm{W}^0$. This is accomplished by initializing the orthogonal matrix $\bm{R}$ as the identity, which aligns the starting point of the fine-tuned model with that of the pretrained parameters. To preserve the orthogonality of $\bm{R}$ throughout optimization, one can employ differentiable orthogonalization techniques introduced by \cite{liu2021orthogonal,lezcano2019cheap}. A detailed discussion follows regarding methods for enforcing orthogonality efficiently.

\subsection{Efficient Orthogonal Parameterization}\label{sect:eff_ortho}

Traditional orthogonalization procedures, such as the differentiable Gram-Schmidt process, are often computationally demanding in practical scenarios~\cite{liu2021orthogonal}. To improve computational tractability, our method leverages the Cayley parameterization to generate orthogonal matrices. Concretely, an orthogonal matrix is constructed as $\thickmuskip=2mu \medmuskip=2mu \bm{R}=(\bm{I}+\bm{Q})(I-\bm{Q})^{-1}$, where $\bm{Q}$ is a skew-symmetric matrix, i.e., $\thickmuskip=2mu \medmuskip=2mu \bm{Q}=-\bm{Q}^\top$. While this strategy is more efficient, it imposes a constraint: only orthogonal matrices with determinant $1$ can be formed, corresponding to the special orthogonal group. Nevertheless, our empirical observations suggest that this constraint does not hinder performance. Despite utilizing the Cayley transform, for large $d$, representing a full $\bm{R}$ remains parameter-intensive. Therefore, to reduce parameter complexity, we express $\bm{R}$ as a block-diagonal matrix composed of $r$ blocks, which results in the following structure:

\begin{equation}
    \footnotesize
    \bm{R}=\text{diag}(\bm{R}_1,\bm{R}_2,\cdots,\bm{R}_r)=\begin{bmatrix}
    \bm{R}_1\in \text{O}(\frac{d}{r}) &  &\\
    &\ddots & \\
    & & \bm{R}_r\in \text{O}(\frac{d}{r})
    \end{bmatrix}\in \text{O}(d)
\end{equation}

Here, $\text{O}(d)$ represents the orthogonal group in $d$ dimensions. The matrices $\bm{R}\in\mathbb{R}^{d\times d}$ and, for every $i$, $\bm{R}_i\in\mathbb{R}^{d/r\times d/r}$ are both orthogonal. If we choose $r=1$, the block-diagonal orthogonal matrix reduces to an unconstrained orthogonal matrix. For an orthogonal matrix of dimensions $d\times d$, the total parameter count is $d(d-1)/2$, leading to a parameter complexity of $\mathcal{O}(d^2)$. In the case of an $r$-block diagonal orthogonal matrix, the parameter count becomes $d(d/r-1)/2$, corresponding to a complexity of $\mathcal{O}(d^2/r)$. Further efficiency can be achieved by reusing a single block across the matrix—formally, $\bm{R}_i=\bm{R}_j$ for all $i\neq j$—which reduces the complexity to $\mathcal{O}(d^2/r^2)$. Importantly, these modifications maintain the orthogonality of $\bm{R}$, thus preserving the hyperspherical energy, and entail no loss in this regard.

We compare OFT and LoRA with respect to parameter efficiency. For LoRA, which uses a low-rank parameter $r'$, the number of trainable parameters totals $r'(d+n)$. Assuming both $r$ and $r'$ scale with the neuron dimension $d$ (for example, setting $r = r' = \alpha d$ where $0 < \alpha \leq 1$ is a fixed constant), LoRA’s parameter count is $\mathcal{O}(d^2 + dn)$, whereas OFT’s becomes $\mathcal{O}(d)$. To illustrate, consider a scenario where the weight matrix is $128\times128$: if LoRA uses $r' = 8$, there are $2,048$ trainable parameters, while OFT with $r = 8$ (and without block sharing) only requires $960$ parameters.

\subsection{Constrained Orthogonal Finetuning}

It is possible to restrict the expressivity of OFT even further by enforcing that the finetuned model remains close to the original pretrained model. This approach, COFT, utilizes the following computation in the forward pass:

\begin{equation}\label{eq:coft_general}
    \footnotesize
    \bm{z}=\bm{W}^\top\bm{x}=(\bm{R}\cdot\bm{W}^0)^\top\bm{x},~~~\text{s.t.}~\bm{R}^\top\bm{R}=\bm{R}\bm{R}^\top=\bm{I},~\norm{\bm{R}-\bm{I}}\leq \epsilon
\end{equation}

This formulation imposes both an orthogonality constraint and a proximity constraint (within $\epsilon$) to the identity matrix. The former can be realized using the Cayley parameterization as described in Section~\ref{sect:eff_ortho}. However, directly embedding the $\epsilon$-deviation constraint within the Cayley-parameterized form is not straightforward. To facilitate analysis, one may apply the Neumann series to approximate $\bm{R}=(\bm{I}+\bm{Q})(\bm{I}-\bm{Q})^{-1}$, yielding $\bm{R}\approx\bm{I}+2\bm{Q}+\mathcal{O}(\bm{Q}^2)$ (assuming that the Neumann {\looseness=-1 \parfillskip=0pt\par}

series converges in the operator norm). Consequently, we are able to express the constraint $\thickmuskip=2mu \medmuskip=2mu\|\bm{R}-\bm{I}\|\leq\epsilon$ directly in terms of the Cayley transform; this is equivalent to enforcing $\thickmuskip=2mu \medmuskip=2mu \|\bm{Q}-\bm{0}\|\leq\epsilon'$, where $\bm{0}$ denotes the zero matrix, and $\epsilon'$ serves as an alternative hyperparameter that may differ from $\epsilon$. Projected gradient descent can be utilized to enforce this revised constraint on the matrix $\bm{Q}$. For identity initialization of the orthogonal matrix $\bm{R}$, $\bm{Q}$ is initialized as the zero matrix. COFT imposes two direct constraints: one that enforces minimal hyperspherical energy variation and another that limits the deviation from the original pretrained model. While the latter constraint is typically assumed implicitly in most finetuning approaches, COFT incorporates it explicitly. Despite the strong performance offered by OFT, our experiments indicate that COFT delivers further improvements in finetuning stability as a result of the explicit deviation regularization. Figure~\ref{fig:eps_coft} illustrates how variations in $\epsilon$ influence COFT's performance. This hyperparameter essentially determines the degree of adaptation permitted during finetuning. As $\epsilon$ increases, models fine-tuned with COFT increasingly resemble those adjusted with OFT. Conversely, smaller $\epsilon$ values restrict adaptation, yielding models that behave similarly to the pretrained text-to-image diffusion model.

\subsection{Re-scaled Orthogonal Finetuning}

We introduce a straightforward augmentation of the basic OFT procedure, wherein we learn an individual scaling parameter for each neuron. The underlying rationale is that scaling the output of neurons does not impact the hyperspherical energy, as this energy is independent of magnitude (which is subsequently normalized). Specifically, the modified forward computation is defined as $\thickmuskip=2mu \medmuskip=2mu\bm{z}=(\bm{R}\bm{W}^0\bm{D})^\top\bm{x}$\footnote{\textbf{Errata}: The NeurIPS camera-ready version contains a typographical error in the re-scaled OFT forward computation, which is written as $\thickmuskip=2mu \medmuskip=2mu\bm{z}=(\bm{D}\bm{R}\bm{W}^0)^\top\bm{x}$. However, the released implementation uses the correct order, so the results reported in Appendix~\ref{app:rescaled} remain valid.}, where $\thickmuskip=2mu \medmuskip=2mu \bm{D}=\text{diag}(s_1,\cdots,s_n)\in\mathbb{R}^{n\times n}$ is a diagonal matrix with trainable positive diagonal entries $s_1,\cdots,s_n$. Unlike the original OFT forward pass in Eq.~\eqref{eq:oft_general}, where only the orthogonal matrix $\bm{R}$ is subject to optimization, here both $\bm{R}$ and the scaling matrix $\bm{D}$ are trained. Incorporating neuron-specific scaling improves the expressiveness of OFT with minimal additional parameter overhead. To isolate the influence of the orthogonal transformation itself, our main experiments employ the original OFT; however, empirical evidence in Appendix~\ref{app:rescaled} shows that the re-scaled variant generally yields even better performance.

\section{Intriguing Insights and Discussions}

\textbf{OFT's applicability across architectures}. OFT is, in theory, compatible with any neural network architecture. In the Transformer setting, LoRA is traditionally applied to the attention matrices~\cite{hulora2022}; for fair comparison, we also constrain OFT to attention weights in our studies. Additionally, OFT is highly compatible with convolutional layers, owing to the interpretable nature of the block-diagonal form of $\bm{R}$ in this context, in contrast to LoRA. If the number of blocks in $\bm{R}$ matches the count of input channels, each block modifies only one specific neuron channel, analogous to learning depth-wise convolutional kernels~\cite{chollet2017xception}. Conversely, when blocks are shared across channels, a single orthogonal matrix transforms all neuron channels uniformly. Preliminary experiments exploring OFT for convolutional layers are reported in Appendix~\ref{app:convolution}.

\textbf{Connection to LoRA}. LoRA introduces a low-rank matrix to its update, thereby safeguarding the integrity of the pretrained weights by avoiding significant modifications. In comparison, OFT restricts the transformation applied to the pretrained weight matrix to be orthogonal (and thus full-rank), which also protects pretraining information by ensuring the transformation cannot obliterate these learned features. The forward computation in OFT can be reformulated as $\thickmuskip=2mu \medmuskip=2mu \bm{z}=(\bm{R}\bm{W}^0)^\top\bm{x}=(\bm{W}^0+(\bm{R}-\bm{I})\bm{W}^0)^\top\bm{x}$, where $\thickmuskip=2mu \medmuskip=2mu (\bm{R}-\bm{I})\bm{W}^0$ serves a similar role to LoRA's low-rank parameter adaptation. Notably, when $\bm{W}^0$ is of full rank, OFT is capable of producing a low-rank update provided $\thickmuskip=2mu \medmuskip=2mu \bm{R}-\bm{I}$ itself has low rank. Both methods introduce hyperparameters for efficiency: LoRA leverages a rank parameter $r'$, while OFT utilizes a block-diagonal parameter $r$ to constrain the parameter count. It is noteworthy that LoRA and OFT reflect two orthogonal philosophies for parameter efficiency—LoRA utilizes low-rank adaptation, whereas OFT capitalizes on block-diagonal sparsity via orthogonality constraints.

\textbf{Why OFT converges faster}. Examination of Figure~\ref{fig:toy_ae} illustrates that modifying neuron orientations is the most efficient strategy for altering semantic content, which aligns with the design of OFT. Additionally, OFT's optimization can be interpreted as adapting neurons on a smooth hyperspherical manifold, resulting in a more favorable optimization surface. Empirical support for this property is reported in \cite{liu2021orthogonal}.

\textbf{Why not minimize hyperspherical energy}. In contrast to \cite{liu2021orthogonal}, our approach does not target the minimization of hyperspherical energy. In classification tasks, redundancy among neurons should be avoided. Achieving minimum hyperspherical energy enforces a uniform distribution of neurons across the hypersphere, but such uniformity is not aligned with the objectives of finetuning, as it risks degrading valuable information acquired during pretraining.

\textbf{Balancing flexibility and regularity during finetuning}. Our findings highlight an inherent tension between flexibility and regularity in finetuning strategies. While standard finetuning offers maximum adaptability, it often compromises stability and may lead to model collapse. Notably, OFT strikes a favorable compromise between these competing needs by maintaining pairwise neuron angular relationships, despite its conceptual simplicity. The block-diagonal approach serves as an even more stringent form of regularization, constraining the orthogonal matrix.

\textbf{Inference efficiency preserved}. In contrast to ControlNet, the OFT approach imposes no extra computational burden during inference on the finetuned model. At inference, we directly integrate the trained orthogonal matrix $\bm{R}$ with the pretrained weight matrix $\bm{W}^0$, resulting in an effective weight matrix $\thickmuskip=2mu \medmuskip=2mu \bm{W}=\bm{R}\bm{W}^0$. Consequently, the inference performance matches that of the original pretrained architecture.

\section{Experiments and Results}

\textbf{Experimental setup}. Our experiments utilize the Stable Diffusion v1.5~\cite{rombach2022high} model for pretrained text-to-image synthesis. For impartial evaluation, image samples are randomly selected across methods. Protocols for subject-driven tasks align with DreamBooth~\cite{ruiz2023dreambooth}, while protocols for controllable generation are consistent with ControlNet~\cite{zhang2023adding} and T2I-Adapter~\cite{mou2023t2i}. To enable fair benchmarking against LoRA, OFT and COFT modifications are restricted to layers also altered by LoRA. Supplementary details and results can be found in Appendix~\ref{app:settings}.

\subsection{Subject-driven Generation}

\textbf{Protocol}. DreamBooth~\cite{ruiz2023dreambooth} and LoRA~\cite{hulora2022} function as baseline comparators. All evaluated techniques employ the same objective function as DreamBooth. Baseline approaches are implemented in accordance with their respective original publications, utilizing optimal hyperparameters. Expanded results are included in Appendix~\ref{app:settings},\ref{app:coft_oft},\ref{app:more_qualitative},\ref{app:failure}.

\textbf{Finetuning stability and convergence}. We begin by analyzing the stability and speed of convergence during finetuning for DreamBooth, LoRA, OFT, and COFT. The corresponding outcomes are illustrated in Figure~\ref{fig:teaser} and Figure~\ref{fig:db_convergence}. It is evident from these figures that both COFT and OFT facilitate notably stable finetuning processes for diffusion models. Within 400 iterations, effective control is attained by DreamBooth and the OFT variants, whereas LoRA struggles to maintain subject identity. By the 2000-iteration mark, DreamBooth's outputs start to degrade with mode collapse, and LoRA does not successfully synthesize a yellow shirt, producing yellow fur instead. In stark contrast, OFT and COFT demonstrate continued robust and coherent control over the generated outputs. These findings confirm the efficient and stable convergence properties of the OFT approach within the context of subject-driven generation. An additional advantage is the low sensitivity of OFT and COFT to the total number of finetuning iterations; this feature simplifies training, eliminating the need for precise tuning of iteration counts for different subjects. It is therefore feasible to select a generously large iteration number for OFT and COFT without risk. Furthermore, COFT with an appropriately chosen $\epsilon$ removes the challenge of configuring both the learning rate and iteration number, simplifying hyperparameter selection.

\textbf{Quantitative comparison}. Adopting the evaluation protocol of \cite{ruiz2023dreambooth}, we carry out quantitative assessments of subject fidelity (using DINO~\cite{caron2021emerging} and CLIP-I~\cite{radford2021learning}), text prompt fidelity (CLIP-T~\cite{radford2021learning}), and sample diversity (LPIPS~\cite{zhang2018unreasonable}). CLIP-I is determined by averaging pairwise cosine similarities of CLIP embeddings derived from both generated and real images, while DINO provides an analogous measure based on ViT S/16 DINO embeddings. The CLIP-T metric is computed as the average cosine similarity between the text prompts and the corresponding generated image embeddings. To measure diversity, we calculate the average LPIPS cosine similarity among images generated for the same subject using identical text prompts. According to Table~\ref{table:dreambooth_final}, both COFT and OFT surpass DreamBooth and LoRA by significant margins on the DINO and CLIP-I measures, and also exhibit slightly improved or comparable outcomes for prompt fidelity and diversity. For robustness, each method is tested by finetuning the same pretrained network across 30 random seeds, thereby reducing variance from stochasticity. Overall, these results indicate that OFT not only ensures improved convergence and training stability, but also delivers consistently higher final accuracy and fidelity.

\textbf{Qualitative comparison}. To provide a clearer illustration of the advantages offered by OFT, Figure~\ref{fig:dreambooth_final} presents a selection of randomly chosen outputs for subject-driven generation. To ensure a just comparison, each example is produced using the same finetuned model with each respective approach, and no text prompt is specifically tailored or re-optimized for a given result. For all methods, we utilize the model checkpoint yielding the highest validation CLIP score. From the visualizations in Figure~\ref{fig:dreambooth_final}, it is evident that OFT and COFT are both able to maintain high fidelity to the intended semantic subject, while LoRA frequently struggles with subject identity retention (for instance, LoRA fails to preserve the bowl’s subject identity in its respective row). Furthermore, OFT and COFT display notably improved accuracy in controlling generation via text prompts, in contrast to DreamBooth which, despite maintaining subject identity, regularly generates outputs that deviate from the provided textual instruction (such as in the initial bowl example). These qualitative results collectively demonstrate that our OFT framework excels in balancing both controllability and faithful subject retention. In addition, OFT exhibits robustness to the choice of iteration count, performing reliably even at higher numbers, whereas DreamBooth and LoRA do not share this trait. Additional qualitative comparisons can be found in Appendix~\ref{app:more_qualitative}. A human study described in Appendix~\ref{app:human} further substantiates OFT’s superior performance.

\subsection{Controllable Generation}

\textbf{Settings}. For baseline comparisons, we include ControlNet~\cite{zhang2023adding}, T2I-Adapter~\cite{mou2023t2i}, and LoRA~\cite{hulora2022}. Our evaluation encompasses three difficult controllable generation tasks: Canny edge-to-image~(C2I) on the COCO dataset~\cite{lin2014microsoft}, segmentation map-to-image~(S2I) on ADE20K~\cite{zhou2017scene}, and landmark-to-face~(L2F) on CelebA-HQ~\cite{karras2018progressive,xia2021tedigan}. All models are finetuned on Stable Diffusion~(SD) v1.5 for 20 epochs across the three datasets. Extended results are reported in Appendix~\ref{app:more_qualitative},\ref{app:more_control_task},\ref{app:failure}.

\textbf{Convergence}. To assess how quickly each method converges, we compare ControlNet, T2I-Adapter, LoRA, and COFT on the L2F task. Our evaluation incorporates both numerical and visual analyses. For quantitative comparison, we utilize the mean $\ell_2$ distance calculated between the control and predicted facial landmarks as our main metric. In Figure~\ref{fig:face_convergence}, we depict how face landmark errors change as each model is finetuned over varying numbers of epochs. The results indicate that COFT and OFT display notably faster convergence rates. For instance, LoRA needs 20 epochs to reach the level of performance that OFT achieves after just 8 epochs. It is worth emphasizing that OFT and COFT maintain a number of trainable parameters on par with LoRA—both being considerably leaner than ControlNet—yet they converge much more rapidly compared to previous methods. Furthermore, OFT’s rapid convergence is corroborated by Figure~\ref{fig:teaser}; the right panel shows that OFT attains strong convergence using only 5\% of the ADE20K dataset, highlighting greater data efficiency over both ControlNet and LoRA. For the qualitative assessment, we limit our comparison to OFT, COFT, and ControlNet, as ControlNet attains landmark errors most similar to ours. Figure~\ref{fig:face_convergence_qual} illustrates that OFT and COFT not only converge reliably but also incrementally improve the alignment between generated face poses and target landmarks. By contrast, ControlNet demonstrates an abrupt onset of controllability after the 8-th epoch—mirroring the sharp convergence patterns noted in \cite{zhang2023adding}. This stable and gradual convergence behavior inherent to OFT substantially facilitates its practical use, yielding more consistent and foreseeable training progress. Consequently, identifying optimal hyperparameters for OFT is a more straightforward process.

\textbf{Quantitative comparison}. To assess controllable generation, we propose a control consistency metric. This metric operates by extracting the control signal from the synthesized image and comparing it against the original input control signal. Specifically, for the C2I task, we measure IoU and F1 score; for the S2I task, we report mean IoU alongside both mean and overall accuracy; and for the L2F task, we use the mean $\ell_2$ distance calculated between control landmarks and predicted landmarks. Additional explanations of these consistency metrics can be found in Appendix~\ref{app:settings}. Each method under comparison is evaluated using optimal hyperparameter configurations. As indicated by the results in Table~\ref{table:control_gen}, OFT and COFT consistently deliver superior and more reliable control compared to alternative methods. Our observations indicate that adapter-based methods (\emph{e.g.}, T2I-Adapter and ControlNet) not only exhibit slower convergence but also produce less favorable outcomes. While LoRA outperforms ControlNet in the S2I setting, it lags behind in the C2I and L2F tasks. In most cases, we notice that LoRA's upper performance limit is comparatively restricted, even when extensively tuned. In contrast, OFT does not appear to have reached its upper bound in performance, as empirical evidence suggests that increasing the number of trainable parameters continues to yield improvements. It is worth highlighting that our quantitative evaluation for controllable generation represents an important innovation of our work, as it enables precise assessment of control effectiveness in finetuned models, conditioned on the accuracy limits of pre-existing segmentation or detection networks.

\textbf{Qualitative comparison}. We conduct a qualitative evaluation of OFT and COFT against leading approaches such as ControlNet, T2I-Adapter, and LoRA. As illustrated by the randomly generated samples in Figure~\ref{fig:control_final}, OFT and COFT consistently deliver images that are not only high in fidelity and realism but also exhibit precise control capabilities. In the S2I setting, LoRA is unable to synthesize images that correspond to the provided segmentation map, whereas ControlNet, OFT, and COFT successfully regulate generation in alignment with input controls. Compared with ControlNet, our OFT and COFT approaches achieve similarly or more detailed and visually plausible images, capturing geometric accuracy with significantly fewer model parameters. For the C2I task, OFT and COFT generate semantically consistent images from basic Canny edge inputs, outperforming T2I-Adapter and LoRA, which struggle with this level of abstraction. In the L2F task, our methods exhibit the most accurate pose alignment, even for faces under difficult orientations. Across these three tasks, our results demonstrate that both OFT and COFT offer superior qualitative outcomes relative to existing state-of-the-art baselines, underscoring the practical strengths of our OFT framework for controllable image generation. Additional qualitative comparisons spanning all three control tasks can be found in Appendix~\ref{app:control_exp}. Furthermore, Appendix~\ref{app:more_control_task} presents evidence that OFT generalizes well to other control applications, such as dense pose to human body, sketch to image, and depth to image translation.

\section{Concluding Remarks and Open Problems}

Inspired by the insight that angular relationships among neurons play a vital role in shaping visual semantics, we introduce a straightforward yet robust finetuning strategy—orthogonal finetuning—for managing text-to-image diffusion models. Our approach addresses two main scenarios: subject-driven generation and controllable generation. Relative to current approaches, OFT offers superior controllability and greater finetuning stability, all while requiring fewer finetuning parameters. Crucially, OFT imposes no extra inference burden, supporting deployment in resource-constrained environments.

The proposed OFT approach also brings to light several intriguing unresolved questions. Firstly, while OFT enforces orthogonality using Cayley parametrization—which necessitates matrix inversion—this operation may hinder the scalability of the method. Although we mitigate this limitation through block diagonal parametrization, developing a more efficient, differentiable technique for matrix inversion remains unresolved. Secondly, OFT presents notable compositional properties: orthogonal matrices generated by multiple OFT finetuning sessions can be sequentially multiplied, resulting in another orthogonal matrix. An open question is whether the composite of these matrices continues to encode all relevant information from the individual downstream tasks. Lastly, the method’s parameter efficiency relies heavily on the adoption of a block diagonal structure, which inherently introduces biases and restricts adaptability. Identifying pathways to enhance parameter efficiency while minimizing bias and maximizing flexibility is an important direction for future research.